\documentclass[11pt,a4paper]{article}
\usepackage[margin=2.4cm]{geometry}
\usepackage{amsmath,amssymb}
\usepackage[colorlinks=true,linkcolor=blue,urlcolor=blue]{hyperref}
\usepackage{xcolor}
\usepackage{parskip}

\newenvironment{reviewercomment}%
{\par\medskip\noindent\begin{quotation}\itshape\color{black!70}}%
{\end{quotation}\medskip}
\newcommand{\response}{\par\noindent\textbf{Response:}\ }
\newcommand{\changes}{\par\noindent\textbf{Changes made:}\ }

\begin{document}

\begin{center}
{\Large\bfseries Response to Reviewer \#1}\\[0.4em]
{\large Ms.\ Ref.\ No.: MEAS-D-26-10707}\\[0.3em]
Gradient-Sensitive Functional Descriptors for Interferometric Surface Metrology:\\
Complementing Amplitude-Averaging ISO Parameters\\[0.3em]
\textit{Measurement} --- Revised version, September 2026
\end{center}

\bigskip

I thank Reviewer~\#1 for the positive assessment --- describing the work as a complete interferometric pipeline and not questioning the correctness of the measurements --- and for three precise, constructive points. Each is addressed below. Reviewer comments are reproduced in italics; my responses give the change and its location in the revised manuscript. A marked (redline) version of the manuscript accompanies this letter.

At the outset I confirm the central reframing the reviewer's comments motivate. I no longer present the work as ``Banach versus ISO''. The defensible thesis, now adopted throughout, is: the $L^1$ and $L^2$ norms establish a functional-analytic representation \emph{consistent with} the established ISO amplitude descriptors ($R_a$, $R_q$), whereas the bounded-variation and Sobolev gradient descriptors provide \emph{complementary} information about spatial variation that amplitude-only descriptors do not capture. Functional analysis is used as an organising framework, not as a competitor to ISO.

\subsection*{Point 1.1 --- Conceptual overclaiming of the $R_a = L^1$, $R_q = L^2$ equivalences}
\begin{reviewercomment}
There is some conceptual overclaiming. The numerical equivalence Ra = L\textsuperscript{1} and Rq = L\textsuperscript{2} is presented as an ``identity'' (which it is, by definition), but the framing suggests more novelty than is justified. Labelling a standard arithmetic mean of absolute deviations as an L\textsuperscript{1} norm provides no new metrological information. The functional-analytic terminology, while mathematically precise, sometimes obscures rather than illuminates.
\end{reviewercomment}
\response I agree and have removed every framing that could suggest novelty in these definitional equalities. The manuscript already stated that the equalities ``are not empirical findings; they are definitional consequences'' and that relabelling $R_a$ as $L^1$ ``provides no new metrological information'' (Sec.~5.3, \emph{Mathematical equivalences}); the revision now applies this stance at every location where the equalities appear.
\changes (a)~In Sec.~4.4 the sentence beginning ``The numerical identity\ldots confirms\ldots'' has been rewritten to state that the equalities ``are definitional --- each pair is the same formula applied to the same centred residual --- and are reported purely as a computational cross-check that both descriptor families are evaluated on identical data, not as an empirical finding.'' (b)~The closing paragraph of the Introduction now states explicitly: ``No claim of novelty is attached to the norms themselves --- they are standard tools of functional analysis, and the functional-analytic vocabulary is used only as an organising taxonomy; the controlled comparison framework is the contribution.'' (c)~The title already frames the descriptors as \emph{complementing} amplitude-averaging ISO parameters, consistent with this position.

\subsection*{Point 1.2 --- Missing base references on surface metrology and fringe analysis}
\begin{reviewercomment}
The introduction does not cover the base and actual researches on surface metrology, e.g., \url{https://doi.org/10.1364/JOSA.72.000156}, \url{https://doi.org/10.1016/J.OPTLASENG.2005.10.012}, \url{https://doi.org/10.3390/asi9020031}, \url{http://dx.doi.org/10.1117/1.AP.1.2.025001}.
\end{reviewercomment}
\response All four references have been added and are discussed individually in a new Introduction paragraph on the methodological lineage of fringe-pattern analysis, which also positions the present radial excess-fraction method within the single-frame, non-phase-shifting family.
\changes New Introduction paragraph discussing, each in its own sentence: Takeda et al.\ (Fourier-transform method of fringe-pattern analysis, J.~Opt.~Soc.~Am.\ 72:156, 1982); Kemao (two-dimensional windowed Fourier transform, Opt.\ Lasers Eng.\ 45:304, 2007); Feng et al.\ (fringe pattern analysis using deep learning, Adv.\ Photonics 1:025001, 2019); and Galaktionov \& Toporovsky (multi-stage fringe-map reconstruction for non-phase-shifting interferometry, Appl.\ Syst.\ Innov.\ 9:31, 2026).

\subsection*{Point 1.3 --- Redundancy of BV density and the $W^{1,2}$ seminorm ($r = 0.94$)}
\begin{reviewercomment}
The correlation between BV density and W\textsuperscript{1,2} seminorm (r=0.94) suggests these descriptors may be providing essentially the same information, yet they are presented as complementary. The paper could more clearly distinguish which descriptors offer truly independent information. The statistical power is limited (n=5 sweeps), and the reported 95\% expanded uncertainties should be treated as lower bounds.
\end{reviewercomment}
\response I agree, and I have both clarified the redundancy and stated precisely where the two descriptors diverge. As \emph{repeatability metrics across sweeps} they are largely redundant ($r=0.94$), and the revised text now says routine reporting of both is not warranted. They are, however, strongly correlated but \emph{not mathematically equivalent}: BV integrates the first power of the local slope, $TV=\int|p'|\,ds$, whereas $W^{1,2}$ integrates its square, $\int|p'|^2\,ds$, so they coincide only when the slope magnitude is nearly constant and diverge when the slope distribution is uneven. The defect responses show exactly this divergence: for the same 100~nm spike, $W^{1,2}$ responded with $+622\%$ versus $+18\%$ for BV density, because quadratic weighting amplifies a single steep event far more strongly than the linear measure. A descriptor-space map is now given: (i)~$R_a$/$R_q$ carry amplitude; (ii)~$R_{sk}$ carries an independent asymmetry axis; (iii)~BV density carries the localised-gradient axis and is recommended as the primary gradient descriptor; (iv)~$W^{1,2}$ adds value mainly as a high-sensitivity alarm for isolated steep features; (v)~$L^4/L^1$ provides a distribution-shape axis. On statistical power: the lower-bound status of all uncertainties is now stated in the Abstract as well as in Sec.~3.7 and the Limitations.
\changes Sec.~5.3 (\emph{Practical considerations}) rewritten with the mathematical distinction and the descriptor-independence map; Abstract updated to state that all uncertainties are within-run lower bounds.

\subsection*{Point 1.4 --- Uncertain practical value; the Flick flat is a known, certified feature}
\begin{reviewercomment}
The practical value remains uncertain. While BV density detects the Flick flat, it is unclear whether this provides actionable information beyond what an operator might see in a polar plot inspection. The paper does not demonstrate an application where gradient-based descriptors reveal a defect that would otherwise be missed in quality control --- the Flick flat is certified, so its presence is already known.
\end{reviewercomment}
\response This is fair, and the revision states it openly. The Flick flat analysis is an existence proof under ground-truth conditions, not a blind detection study. The narrower practical argument I now make is that the gradient-sensitive layer converts what an operator \emph{might} notice in a visual polar-plot inspection into a quantitative, thresholdable, automatable indicator computed from data already acquired for ISO reporting --- with no operator in the loop, a defined angular-position estimate, and a repeatability statement. This is precisely the answer to ``why BV if I already have $R_a$?'': $R_a$ is an amplitude average, whereas BV responds to localised profile variation. I explicitly acknowledge that a demonstration in which gradient descriptors reveal a defect that routine quality control would otherwise miss is still outstanding.
\changes New paragraph \emph{Practical value and the known-feature caveat} in Sec.~5.3; Future Work item~(iii) now specifies a blind-detection protocol (analyst unaware of defect presence and position, artefacts with independently characterised defects of calibrated dimensions) as the decisive next test of practical value; the corresponding Limitations entry names blind-detection performance as uncharacterised.

\bigskip
I hope these changes address the reviewer's points fully.

\bigskip
Yours sincerely,\\[0.6em]
Dawid Kucharski\\
Institute of Mechanical Technology, Poznan University of Technology\\
\href{mailto:dawid.kucharski@put.poznan.pl}{dawid.kucharski@put.poznan.pl}

\end{document}
